\chapter{Method}
\section{Experimental Setup}
General notes on the agility and openness of the methods applied in this study can be found in the \hyperref[sec:appendixA]{Appendix A}. Now, it is time to have a closer look at the separate steps in the workflow:
\begin{enumerate}
         \item Create a dataset.
         \begin{enumerate}
		\item Specify the text corpus to be used.
		 \begin{enumerate}
			\item Segment the corpus into smaller pieces, e.g. sentences.
			\item Tokenize every segment.
		\end{enumerate}
		\item Specify a gold standard for hyponymy.
		\item Define hyponymy pairs for every segment in the corpus.
	\end{enumerate}
         \item Split the dataset into subsets for training, validation and evaluation.
         \item Implement a transformer language model \parencite{vaswaniAttentionAllYou2017} for hyponymy extraction.
         \begin{enumerate}
		\item Train the model on the training subset.
		\item Use the validation subset to make sure that the model inferred general rules instead of memorizing input patterns.
	\end{enumerate}
         \item Evaluate the model using the evaluation dataset.
\end{enumerate}
Since machine learning is the major approach to be considered here, the size of the dataset (especially the training subset) is crucial \parencite[13]{hestnessDeepLearningScaling2017}. The rule of maximizing training data applies not only to machine learning, but also to other methods of distributional semantics, e.g. in the case of semantic relations \parencite[27]{herbelotBuildingSharedWorld2015}. Therefore, it seems reasonable to prefer the Corpus Corporum (about 190 million tokens of plain text, cf. \textcite{roelliCorpusCorporumNew2014}) over PROIEL (about 220.000 tokens of annotated text). Unfortunately, this limits the usage of linguistic research results (i.e. annotations) as input for the language models. On the other side, the practical application of the model for language learning purposes becomes easier and more flexible, in the sense that, theoretically, end users can supply any plain Latin text of their choice, without the need to provide additional annotations.

Before we can start to detect hyponyms, we need to induce word senses, especially for polysemous words. To make sure that fine-grained distinctions between different usages of the same word are recognized by the model, it should be evaluated on a test dataset that has the following structure:
\begin{itemize}
	\item For each segment in a pair of segments (e.g. sentences),
	\item look at the target word(s) and
	\item predict whether it was used in the same sense as in the other segment.
\end{itemize}
The second step can be formalized by providing character indices referring to the position of the target word in a string representation of the segment. This works even for cases where a target word form occurs multiple times in the segment and we want to induce the word sense for just a single specific instance. The third step, i.e. prediction of semantic similarity, can be binary in this very basic operationalization: The two target words are either used in the same sense or not. This dichotomy consciously abstracts from more complicated evidence, e.g. where two instances are used in a very similar, but not quite the same sense. In a subsequent effort, we may try to associate the current instance with other segments where the word usage is similar, thus clustering occurrences of the same word sense together. 

Since there are no semantically annotated authentic text corpora for Latin, the next best thing to use are grammar books and dictionaries: They often distinguish between various word senses and offer exemplary contexts for each one \parencites{georgesAusfuhrlichesLateinischDeutschesHandworterbuch1913}{niederauNavigiumLateinDeutschWorterbuch2012}. Some of them cite entire text passages \parencites{shortLatinDictionary1879}{kuhnerAusfuhrlicheGrammatikLateinischen1914}{mengeLehrbuchLateinischenSyntax2009} to support their classifications. Such references can be seen as annotations: If a dictionary lists various text passages as containing a specific word in the same sense, this is roughly equivalent to semantically annotated datasets where sense identifiers from static sense inventories are assigned to tokens in a text corpus. Now, we could proceed by looking at the largest entries in a dictionary, assuming that they contain the largest number of different word senses. This, however, would amount to accepting the dictionary editor's intuition, which is not a usage-based approach. Instead, we will find the most polysemous words in a corpus inductively, by looking at their usage contexts. In this view, the most polysemous word will be that which is used in the most diverse contexts.

\label{functionWordsZipfDistribution}
A naive formalization of diverse contexts is to count the number of types that occur in the same sentence as a target type in the corpus. This disregards the absolute frequencies of the types in a corpus and will lead to a Zipfian distribution \parencite[VI]{zipfPsychobiologyLanguageIntroduction1936} of target types. Thus, function words like \textit{et, in, ut, est, non, cum} `and, in, as, is, not, with' will dominate the ranking just because of their immense absolute frequency in the corpus. We can limit their impact by randomly down-sampling their occurrences using an arbitrary upper bound of 100.
\label{high_frequency_hyponyms_for_language_learning}
Since our goal is to produce exercises for language learning, the potential hyponyms should be rather frequent in a given corpus \parencites[27]{ellisUsagebasedLanguageInvestigating2013}[2]{robillardMonolingualBilingualChildren2014}. It is therefore reasonable to establish a minimum threshold for word frequency. This criterion interacts with Latin being a highly inflectional language, thus producing many morphologically different forms for the same lexeme (see Table \ref{tableTypesPerLemma}).
\begin{table}[ht]
	\begin{tabular}{ c | c | c | c }
  		 & PROIEL & French Treebank & Hamburg Dependency Treebank \\ \hline
  		Types per Lemma & 3.64 & 1.21 & 2.09 \\
	\end{tabular}
	\caption{Average number of types per lemma for Latin, French and German.}
	\label{tableTypesPerLemma}
\end{table}

%TODO: did we really use lemmatized texts?
Thus, we will perform this initial experiment on professionally lemmatized texts, which restricts us to the PROIEL treebank. Using lemmata instead of types has the advantage of increased frequency counts per item (esp. in small corpora) and allows us to easily group contexts by lexeme. The downside here is that the diversity of contexts cannot be distinguished anymore for various word forms of the same lexeme, but we can reintroduce that distinction in a later step. The comparison of two contexts for a given target word will be made using representations of sentences in a \gls{CVS}. This mapping from textual to numerical entities (i.e. vectors) is not straightforward, because it is unclear how exactly we should model the difference between short and long sentences that share roughly the same content, but with different wording. 

%TODO: add comparison of different vector measures, including Euclidean distance, cosine distance etc.
If we pad all shorter sentences to an arbitrary upper bound (e.g. the length of the longest sentence in the corpus) using placeholder values, the model will be able to calculate the cosine similarity\footnote{Throughout this thesis, the terms \textit{cosine similarity} and \textit{cosine distance} will be used interchangeably. The similarity is usually a value between -1 and 1. The distance is usually calculated by subtracting the similarity from 1.} \parencite[3]{luoCosineNormalizationUsing2017} between two sentence vectors as input. However, naively using a fixed value as the placeholder in right-padded sequences \parencite[386]{sachanEffectiveUseBidirectional2018} is equivalent to assuming that, compared to very long sentences, every short sentence tends to repeat the same word over and over again after a certain time. From a linguistic perspective, this formalization seems rather unintuitive because it does not take into account the effects of text length on lexis \parencites[31]{golcherStylometryInterplayTopic2011}[141]{ochabStylometryLiteraryPapyri2019}, which may lead to a clustering of input sentences roughly by (original) sequence length.

\section{Sensitivity to Lexical and Syntactic Variation}
%TODO: describe the transformer model (layers, parameters etc.)
%TODO: explain why positional encoding is important in Latin, even for distributional semantics
Against this backdrop, normalization of text length is probably just as problematic as the disruptive influence of text length in the first place. Possible alternatives consist in manual ad hoc annotation or trying to circumvent the problem. We will proceed by testing a model of semantic diversity that includes lexical and syntactic cues while also ignoring text length. By training a transformer language model \parencite[9]{vaswaniAttentionAllYou2017} on the raw text of the PROIEL corpus, we build a \gls{CVS} into which single tokens can be embedded. Our implementation uses Positional Encoding \parencite[4]{conneauCrosslingualLanguageModel2019} to explicitly analyze every word's position in a given sequence, thereby including basic syntactic notions in the model. Besides, the aforementioned distributional approach makes the representations lexically predictable, i.e. components of collocations will be represented next to each other in the \gls{CVS}. To corroborate these hypotheses, we evaluate our model on a few arbitrarily selected example sentences from PROIEL that are lexically identical, but differ in their word order (see Table \ref{tableCosSimWordOrder}). As a cross check, we have calculated the cosine similarity between entirely identical sentences and verified that it equals 1, which is the maximum value.
\begin{table}[ht]
	\begin{tabular}{p{1.1\linewidth} | p{.06\linewidth}}
  		Context & Cosine Similarity \\ \hline
        \ex[exno=PROIEL 12966]
            	\begingl
            		\gla si \textbf{fu-eri-t} oculus tu-us ... //
            		\glb if be.{\sc pf-sbjv-3sg} eye your-{\sc m} //
            		\glft if your eye was ... //
            	\endgl
        \xe & \\
	\ex[exno=PROIEL 17274]
            	\begingl
            		\gla si oculus tu-us \textbf{fu-eri-t} ... //
            		\glb if eye your-{\sc m} be.{\sc pf-sbjv-3sg} //
            		\glft if your eye was ... //
            	\endgl
        \xe & 0.81 \\ \hline
	\ex[exno=PROIEL 13425]
            	\begingl
            		\gla \textbf{qui} autem non habe-t ... //
            		\glb who but not have-{\sc 3sg} //
            		\glft but who does not have ... //
            	\endgl
        \xe
	& \\
	\ex[exno=PROIEL 17880]
            	\begingl
            		\gla ab e-o autem \textbf{qui} non habe-t ... //
            		\glb from this-{\sc sg.abl.m/n} but who not have-{\sc 3sg} //
            		\glft but from him who does not have ... //
            	\endgl
        \xe
	& 0.76 \\ \hline
	\ex[exno=PROIEL 13883]
		\par
            	\begingl
            		\gla era-t enim habe-ns \textbf{mult-as} possession-es //
            		\glb be.{\sc impf-3sg} for have-{\sc part.pres} many-{\sc f.pl} possession-{\sc pl.nom/acc} //
            		\glft for he had many possessions //
            	\endgl
        \xe
        & \\
	\ex[exno=PROIEL 11010]
		\par
            	\begingl
            		\gla era-t enim habe-ns possession-es \textbf{mult-as} //
            		\glb be.{\sc impf-3sg} for have-{\sc part.pres} possession-{\sc pl.nom/acc} many-{\sc f.pl} //
            		\glft for he had many possessions //
            	\endgl
        \xe
	& 0.37 \\ \hline
	\ex[exno=PROIEL 14258]
		\par
            	\begingl
            		\gla ... praeteri-bi-t \textbf{haec} generatio donec omn-ia haec fi-a-nt //
            		\glb ... pass.by-{\sc fut-3sg} this.{\sc f.sg / n.pl} generation until all-{\sc n.pl} this.{\sc f.sg / n.pl} happen-{\sc sbjv-3pl} //
            		\glft ... this generation will pass by before all this happens //
            	\endgl
        \xe
	& \\
	\ex[exno=PROIEL 18059]
		\par
            	\begingl
            		\gla ... praeteri-bi-t generatio \textbf{haec} donec omn-ia fi-a-nt //
            		\glb ... pass.by-{\sc fut-3sg} generation this.{\sc f.sg / n.pl} until all-{\sc n.pl} happen-{\sc sbjv-3pl} //
            		\glft ... this generation will pass by before everything happens //
            	\endgl
        \xe
	& 0.83 \\ \hline
	\ex[exno=PROIEL 14376]
            	\begingl
            		\gla opus \textbf{bon-um} opera-t-a es-t in me //
            		\glb deed good-{\sc n} do-{part.pf-f} be-{\sc 3sg} towards I.{\sc acc} //
            		\glft she has done a good deed towards me //
            	\endgl
        \xe
	& \\
	\ex[exno=PROIEL 11310]
            	\begingl
            		\gla \textbf{bon-um} opus opera-t-a est in me //
            		\glb good-{\sc n} deed do-{part.pf-f} {\sc aux} towards I.{\sc acc} //
            		\glft she has done a good deed towards me //
            	\endgl
        \xe
	& 0.75 \\ \hline
	\ex[exno=PROIEL 11311]
            	\begingl
            		\gla semper enim \textbf{pauper-es} habe-tis vo-bis cum //
            		\glb always for poor-{\sc pl.f/m} have-{\sc 2pl} you-{abl} with //
            		\glft for you always have poor people with you //
            	\endgl
        \xe
	& \\
	\ex[exno=PROIEL 19419]
            	\begingl
            		\gla \textbf{pauper-es} enim semper habe-tis vo-bis cum //
            		\glb poor-{\sc pl.f/m} for always have-{\sc 2pl} you-{abl} with //
            		\glft for you always have poor people with you //
            	\endgl
        \xe
	& 0.81 \\ \hline
	\end{tabular}
	\caption{Cosine similarity for tokens in almost identical sentences from PROIEL. Target tokens are highlighted in bold.}
	\label{tableCosSimWordOrder}
\end{table}

The outlier value for the sentences 13883 and 11010 indicates a strong change in the token's representation depending on its position in this context. Morphologically, there is an unambiguous agreement between the noun \textit{possessiones} `possessions' and its modifying adjective \textit{multas} `many' in case, number and gender. Thus, further research is to be done here to explain this outcome by testing for other possible interferences, e.g. sentence length, part of speech or lexeme-specific syntactic preferences. It seems that, at least for the PROIEL corpus, the relatively free word order does not automatically imply a semantic equivalence of various syntactic configurations \parencite[452]{devineLatinWordOrder2006}, even though some of the previous scholarly literature suggested this to be the case \parencite[421]{niemeyerZurStellungAttributiven1997}. 

Lexis, on the other hand, is a much more obvious indicator of semantic differences, especially from a distributional perspective. To compare sentences with identical syntax, we arbitrarily pick examples from PROIEL that are longer than 3 tokens and differ in exactly 1 word, which was replaced by either an entirely different word or by a morphological variant of the same lexeme. Such syntactically almost identical sentences can be compared by computing the cosine similarity of their vector representations. These vectors are aggregated by averaging the vectors for all tokens in the sentence, which is a common strategy \parencite[1]{adiFinegrainedAnalysisSentence2016}. Again, as a cross check, we make sure that the cosine similarity for completely identical input sentences is almost at the maximum value. Then, we proceed to compare each sentence to its syntactical twin where 1 token was replaced (see Table \ref{tableCosSimWordSubstitution}).
%TODO: Synonymie zwischen initium und principium genauer erklären >> syntaktische Unterschiede?
\begin{table}[ht]
	\begin{tabular}{p{1.1\linewidth} | p{.06\linewidth}}
  		Context & Cosine Similarity \\ \hline
	\ex[exno=PROIEL 34155]
            	\begingl
            		\gla ego sum Alpha et Omega \textbf{initium} et finis //
            		\glb I be.{\sc 1sg} Alpha and Omega beginning and end //
            		\glft I am Alpha and Omega, \textbf{beginning} and end //
            	\endgl
        \xe & \\
	\ex[exno=PROIEL 33474]
            	\begingl
            		\gla ego sum Alpha et Omega \textbf{principium} et finis //
            		\glb I be.{\sc 1sg} Alpha and Omega origin and end //
            		\glft I am Alpha and Omega, \textbf{origin} and end //
            	\endgl
        \xe
	& 0.90 \\ \hline
	\ex[exno=PROIEL 33726]
            	\begingl
            		\gla et \textbf{secundus} angelus tub-a cecini-t //
            		\glb and second angel trumpet-{\sc abl} sound.pf-{\sc 3sg} //
            		\glft and the \textbf{second} angel sounded his trumpet //
            	\endgl
        \xe & \\
	\ex[exno=PROIEL 33823]
            	\begingl
            		\gla et \textbf{septimus} angelus tub-a cecini-t //
            		\glb and seventh angel trumpet-{\sc abl} sound.pf-{\sc 3sg} //
            		\glft and the \textbf{seventh} angel sounded his trumpet //
            	\endgl
        \xe & 0.88 \\ \hline
	\ex[exno=PROIEL 33661]
            	\begingl
            		\gla et ecce equus \textbf{pallidus} //
            		\glb and there horse pale //
            		\glft and, see, there is a \textbf{pale} horse //
            	\endgl
        \xe & \\
	\ex[exno=PROIEL 33655]
            	\begingl
            		\gla et ecce equus \textbf{niger} //
            		\glb and there horse black //
            		\glft and, see, there is a \textbf{black} horse //
            	\endgl
        \xe & 0.85 \\ \hline
	\ex[exno=PROIEL 14139]
            	\begingl
            		\gla dilig-e-s proxim-um tu-um \textbf{sicut} te ips-um //
            		\glb love-{\sc fut-2sg} neighbor-{\sc acc} your-{\sc acc} like you.{\sc acc} self-{\sc acc} //
            		\glft you shall love your neighbor \textbf{like} yourself //
            	\endgl
        \xe & \\
	\ex[exno=PROIEL 11217]
            	\begingl
            		\gla dilig-e-s proxim-um tu-um \textbf{tamquam} te ips-um //
            		\glb love-{\sc fut-2sg} neighbor-{\sc acc} your-{\sc acc} as you.{\sc acc} self-{\sc acc} //
            		\glft you shall love your neighbor \textbf{as} yourself //
            	\endgl
        \xe & 0.92 \\ \hline
	\ex[exno=PROIEL 20141]
            	\begingl
            		\gla si sic e-um vol-o manere donec \textbf{veni-am} quid ad te //
            		\glb if so this-{\sc m.acc} want-{\sc 1sg} stay until come-{\sc sbjv.1sg} what to you.{\sc acc} //
            		\glft if I want him to stay this way until I \textbf{come}, what is it to you //
            	\endgl
        \xe & \\
	\ex[exno=PROIEL 63139]
            	\begingl
            		\gla si sic e-um vol-o manere donec \textbf{veni-o} quid ad te //
            		\glb  if so this-{\sc m.acc} want-{\sc 1sg} stay until come-{\sc 1sg} what to you.{\sc acc} //
            		\glft if I want him to stay this way until I \textbf{come}, what is it to you //
            	\endgl
        \xe & 0.47 \\ \hline
	\ex[exno=PROIEL 19306]
            	\begingl
            		\gla quaere-ba-nt ergo e-um \textbf{prendere} //
            		\glb seek-{\sc impf-3pl} therefore this-{\sc m.acc} catch //
            		\glft therefore, they wanted to \textbf{get a hold of} him //
            	\endgl
        \xe & \\
	\ex[exno=PROIEL 18951]
            	\begingl
            		\gla quaere-ba-nt ergo e-um \textbf{adprehendere} //
            		\glb seek-{\sc impf-3pl} therefore this-{\sc m.acc} seize //
            		\glft therefore, they wanted to \textbf{seize} him //
            	\endgl
        \xe & 0.81 \\ \hline
	\ex[exno=PROIEL 18848]
            	\begingl
            		\gla qui cred-it in \textbf{me} habe-t vita-m aetern-am //
            		\glb who believe-{\sc 3sg} in I-{\sc acc} have-{\sc 3sg} life-{\sc acc} eternal-{\sc f.acc} //
            		\glft who believes in \textbf{me} has an eternal life //
            	\endgl
        \xe & \\
	\ex[exno=PROIEL 48253]
            	\begingl
            		\gla qui cred-it in \textbf{Fili-um} habe-t vita-m aetern-am //
            		\glb who believe-{\sc 3sg} in Son-{\sc acc} have-{\sc 3sg} life-{\sc acc} eternal-{\sc f.acc} //
            		\glft who believes in \textbf{the Son} has an eternal life //
            	\endgl
        \xe & 0.95 \\ \hline
	\end{tabular}
	\caption{Cosine similarity for sentences from PROIEL in which exactly one token was replaced. Replacements and their translations are highlighted in bold.}
	\label{tableCosSimWordSubstitution}
\end{table}

A major problem with the underlying \gls{CVS} is that its dimensions are not interpretable. In language modeling with neural networks, it is generally assumed that the vector representations entail all kinds of linguistic information, e.g. on morphology, syntax or semantics \parencites[8]{gladkovaAnalogybasedDetectionMorphological2016}[143]{rogersTooManyProblems2017}. However, in each of the linguistic domains, fine-grained nuances may not have been captured adequately, such as the distinction of synonyms from other semantically related words \parencite[115]{karanDistributionalSemanticsApproach2012} or semantic analogies beyond single cases and categories \parencite[142f.]{rogersTooManyProblems2017}. Therefore, it remains unclear why some substitutions invoke much greater changes in similarity than others, e.g. the change of mood (sentences 20141 and 63139) vs. the change of speaker and perspective (sentences 18848 and 48253). An intuitive explanation relates to homographs, i.e. \textit{veniam} could also be an accusative noun (`mercy') instead of a subjunctive verb (`I might come'). Further research in this direction would need to consider alternative sentence representations, inspection of the different attention head layers for the replaced tokens, and systematic trials for specific replacement strategies (inflectional variants, synonyms etc.). For our purposes, it is enough to have a very basic empirical indication that the language model is sensitive to subtle changes in syntax and lexis, which is an important prerequisite for our evaluation of polysemy.

Besides, we may note that most of the replaced tokens seem to exhibit some kind of semantic relation, as can be seen in \gls{LOD} collections like WordNet:
\begin{itemize}
	\item synonymy
	\begin{itemize}
		\item \textit{initium} `beginning' and \textit{principium} `origin' both belong to the same synset (\url{https://latinwordnet.exeter.ac.uk/api/synsets/n/04515071/lemmas/}), defined as `the first part or section of something'.
		\item \textit{adprehendere} `to seize' and \textit{prehendere} `to lay hold of' both belong to the synset `take hold of so as to seize or restrain or stop the motion of' (\url{https://latinwordnet.exeter.ac.uk/api/synsets/v/00986697/lemmas/}).
	\end{itemize}
	\item co-hyponymy
	\begin{itemize}
		\item \textit{secundus} `second' and \textit{septimus} `seventh' are most certainly co-hyponyms in contexts where entities are being enumerated. The appropriate Wordnet synset (\url{https://latinwordnet.exeter.ac.uk/api/synsets/a/02103770/lemmas/}) contains \textit{septimus}, \textit{primus} `first' and \textit{quintus} `fifth' and is defined as `being or denoting a numerical order in a series'. The fact that `secundus' is not yet included there is probably due to the poor quality of the Latin WordNet \parencite[4]{franziniNuncEstAestimandum2019}. Note that there is no direct Latin equivalent for `umpteenth' or other words denoting an arbitrary numerical sequence.
		\item \textit{pallidus} `pale' and \textit{niger} `black' are most certainly co-hyponyms, as they belong to similarly defined synsets: `(used of color) having a dark hue' (\url{https://latinwordnet.exeter.ac.uk/api/synsets/a/00380983/}) and `(used of color) having a relatively small amount of coloring agent' (\url{https://latinwordnet.exeter.ac.uk/api/synsets/a/00380299/}).
		\item \textit{sicut} `so as' and \textit{tamquam} `as much as' both belong to the same synset: `(often followed by \"as\") to the same degree' (\url{https://latinwordnet.exeter.ac.uk/api/synsets/r/00020994/lemmas/}).
	\end{itemize}
\end{itemize}
We may conclude that lexical substitution without syntactic alteration often correlates with semantic relations between the two interchangeable lexemes. In contrast, the remaining two examples from Table \ref{tableCosSimWordSubstitution} exhibit multiple inflections of the same lemma (\textit{venire} `to come'), and a coreference pointing to the figure of Jesus (\textit{me} `me' and \textit{Filium} `Son') in the New Testament. In those cases, none of the traditional semantic relations is present. However, this approach in general is just an expert's intuition and has not been verified empirically by observing actual language use. We therefore need to look for further evidence in the text corpus.

Finding empirical evidence for co-hyponymy is rather difficult in small text corpora, especially when adhering to the Distributional Inclusion Hypothesis \hyperref[distributionalInclusionHypothesis]{mentioned above}, which demands that the hyponyms and their hypernym share a large portion of their usage contexts \parencite[110]{geffetDistributionalInclusionHypotheses2005}. This ambitious requirement is often approximated using more readily available measures, such as graph density in distributional thesauri \parencites[2]{janaNetworkFeaturesBased2018}[4]{janaUsingDistributionalThesaurus2020}.

Building on our heuristic approach from above, we may fine-tune the procedure by considering only the cases where at least one of the two replaced words is a noun, and by excluding all sentence pairs with mere morphological inflection instead of lexical substitution (see Table \ref{tableCosSimWordSubstitutionNouns}). 

\begin{table}[ht]
	\begin{tabular}{p{\linewidth}}
    \ex[exno=PROIEL 60041]
            \begingl
                \gla sext-a \textbf{feri-a} in Syon //
                \glb sixth-{\sc f.sg.nom} holiday-{\sc f.sg.nom} to Syon //
                \glft on the sixth holiday [they move] to Syon //
            \endgl
    \xe \\
	\ex[exno=PROIEL 64252]
            	\begingl
            		\gla sext-a \textbf{di-e} in Syon //
            		\glb sixth-{\sc f.sg.abl} day-{\sc f.sg.abl} to Syon //
            		\glft on the sixth day [they move] to Syon //
            	\endgl
    \xe \\ \hline
	\ex[exno=PROIEL 12797]
            	\begingl
            		\gla et ecce vox de \textbf{cael-is} dice-ns //
            		\glb and there voice from heaven-{\sc n.pl.abl} say-{\sc part.pres} //
            		\glft and there a voice from the sky was saying //
            	\endgl
    \xe \\
	\ex[exno=PROIEL 13725]
            	\begingl
            		\gla et ecce vox de \textbf{nub-e} dice-ns //
            		\glb and there voice from cloud-{\sc abl} say-{\sc part.pres} //
            		\glft and there a voice from the cloud was saying //
            	\endgl
    \xe \\ \hline
    \ex[exno=PROIEL 12929]
            	\begingl
            		\gla nonne et \textbf{publican-i} h-oc faci-unt //
            		\glb not also tax.gatherer-{\sc pl} this-{\sc n.acc} do-{\sc 3pl} //
            		\glft do not tax gatherers also do this? //
            	\endgl
    \xe \\
	\ex[exno=PROIEL 12931]
            	\begingl
            		\gla nonne et \textbf{ethnic-i} h-oc faci-unt //
            		\glb not also pagan-{\sc pl} this-{\sc n.acc} do-{\sc 3pl} //
            		\glft do not pagans also do this? //
            	\endgl
    \xe \\ \hline
    \ex[exno=PROIEL 13048]
            	\begingl
            		\gla et a-it ill-i \textbf{Iesus} //
            		\glb and say-{\sc pf.3sg} that-{\sc m.dat} Jesus //
            		\glft and Jesus told him //
            	\endgl
    \xe \\
	\ex[exno=PROIEL 11219]
            	\begingl
            		\gla et a-it ill-i \textbf{scriba} //
            		\glb and say-{\sc pf.3sg} that-{\sc m.dat} scribe //
            		\glft and the scribe told him //
            	\endgl
    \xe \\ \hline
    \ex[exno=PROIEL 13055]
            	\begingl
            		\gla et responde-ns \textbf{centurio} a-it //
            		\glb and answer-{\sc part.pres} commander say-{\sc pf.3sg} //
            		\glft and as a response, the commander said //
            	\endgl
    \xe \\
	\ex[exno=PROIEL 10471]
    	\begingl
        	\gla et responde-ns \textbf{e-is} a-it //
        	\glb and answer-{\sc part.pres} that-{\sc m.pl.dat} say-{\sc pf.3sg} //
          	\glft and responding to them, he said //	
        \endgl
    \xe \\ \hline
    \ex[exno=PROIEL 47503]
    	\begingl
        	\gla vulp-es fovea-s habe-nt et volucr-es cael-i \textbf{tabernacul-a} //
        	\glb fox-{\sc pl} pit-{\sc pl.acc} have-{\sc 3pl} and bird-{\sc pl} sky-{\sc gen} tent-{\sc pl.acc} //
          	\glft the foxes have pits and the birds in the sky have tents //	
        \endgl
    \xe \\
    \ex[exno=PROIEL 62941]
    	\begingl
        	\gla vulp-es fovea-s habe-nt et volucr-es cael-i \textbf{nid-os} //
        	\glb fox-{\sc pl} pit-{\sc pl.acc} have-{\sc 3pl} and bird-{\sc pl} sky-{\sc gen} nest-{\sc pl.acc} //
          	\glft the foxes have pits and the birds in the sky have nests //	
        \endgl
    \xe
	\end{tabular}
	\caption{Nouns in almost identical sentences from PROIEL. Target tokens are highlighted in bold.}
	\label{tableCosSimWordSubstitutionNouns}
\end{table}

From these examples, we can see how the constraints lead to a higher precision, i.e. out of all results found, more and more tend to be instances of hyponymy or co-hyponymy. However, the recall (How many instances of hyponymy did we miss?) is unknown and can only be determined using gold-standard annotations. We will investigate that problem later in chapter XXX. 
%TODO: determine recall for hyponymy detection
Meanwhile, there are other problems to address. The first example in table \ref{tableCosSimWordSubstitutionNouns} can be regarded as co-hyponymous: \textit{feria} `holiday' and \textit{dies} `day' both share a common transitive relation to the same WordNet synset (`a day reckoned from midnight to midnight'\footnote{\url{https://latinwordnet.exeter.ac.uk/api/synsets/n/10878116/}.}). In our case, this interpretation is questionable and should be more strict: The hypernymy chain leads from \textit{feria} (a holiday) to \textit{dies} (any day from sunrise to sunset) and from there to the higher-level synset of `a day reckoned from midnight to midnight'.

The second example is an instance of metonymy or, functionally speaking, \textit{pars pro toto} \parencite[44]{gardenforsSemanticTransformations2015}: \textit{nubes} `cloud' and \textit{caelum} `sky' are both meant to refer to the same phenomenon, namely the sky above from where a voice is ringing out. In that sense, the two words are meronyms (or even almost synonyms) rather than co-hyponyms, let alone strict hyponyms. There is no simple way for separating their relation from that in the first example, but neural networks may still be able to grasp various cues that, when considered in combination \parencite[1]{hintonImprovingNeuralNetworks2012}, enable us to clearly distinguish the two cases.

The \textit{publicani} `tax gatherers' and \textit{ethnici} `pagans' in the third example lead us to an important decision that we have not made yet: Since hyponymy is a transitive relation (see \hyperref[transitivity]{above}), we may consider a common hypernym for both, such as \textit{homo} `human being`. From that perspective, they would act as co-hyponyms. However, the semantic distance from \textit{publicani} `tax gatherers' to \textit{homo} `human being' is much larger than to more closely related hypernyms such as \textit{magistratus} `public functionary'. For \textit{ethnici} `pagans', on the other side, there does not exist a closer hypernym than \textit{homo} `human being' in the Latin WordNet. While the latter aspect remains to be examined empirically, the conceptual point made here, i.e. the difference in semantic distance to a certain common hypernym, reveals a general problem in applied linguistic research on semantic relations: Does the common hypernym need to be at the same semantic distance for both co-hyponyms? In other words: Many semantic relations are highly underspecified when it comes to their operationalization \parencite[5]{baroniHowWeBLESSed2011}. Therefore, we need to establish both a precise definition of semantic distance and a general threshold for that distance when detecting co-hyponymy.

Since we follow, in this work, a distributional semantic approach, the distance in hyponymy relations can only refer to the distributional similarity of the considered words, i.e. hyponymy pairs with a more similar distribution of typical context words will have less semantic distance than those with rather different contexts. Regarding a threshold, there is no obvious intuition where to put it, not even on a discrete scale like the distinct WordNet synsets. There, the two extremes would be 1 (hypernyms must be directly related) or 0 (any amount of steps along the transitivity chain is allowed). Applying that simple binary choice to our continuous distributional perspective, there is no obvious numeric value (such as cosine distance in a \gls{CVS}) that would correspond to the extreme of 1 in WordNet because a cosine distance of 0 corresponds to proportional vectors, which is generally only the case if all usage contexts of two words in a corpus are exactly similar. Therefore, as an initial working hypothesis, we adopt the other extreme and argue that all words with a common hypernym, independent of the respective semantic distance, can be regarded as co-hyponyms. In that sense, \textit{Iesus} `Jesus' and \textit{scriba} `scribe' from the fourth example are distantly related co-hyponyms, their relation being somewhat similar to that in the third example.

Operating at the interface of semantics and syntax, we also need to consider coreference and its resolution. While a token like \textit{eis} `them' could, in other contexts, be read as a neuter form indicating an antecedent like \textit{verba} `words'\footnote{\url{https://cts.perseids.org/api/cts/?request=GetPassage&urn=urn:cts:latinLit:phi0474.phi040.perseus-lat1:21-40}. For ancient text passages that are not present in any of the treebanks, but do exist in the Perseus Digital Library \parencite{cranePerseusDigitalLibrary2023}, we shall use persistent identifiers according to the Canonical Text Services \parencite{clericePerseidsCanonicalText} protocol as part of the CITE architecture \parencite{blackwellCanonicalTextServices2015}.}, it does point to the animate collective of \textit{mater} `mother' and \textit{fratres} `brothers' (PROIEL 10470) in the present case, thus drawing a possible parallel to \textit{centurio} `commander' (PROIEL 13055), which designates an animate entity as well. Insofar, coreference resolution can help us identify potential further semantic relations between nouns. However, the syntactical (indirect object vs. nominal subject) and morphological (dative plural vs. nominative singular) differences between the two considered instances are important contraindications of a close semantic relation. Thus, we need to add syntactic constraints to our model of hyponymy, where a simple and strict starting point might be syntactic equivalence on the token level. Again, this will probably increase the precision of our method, at the expense of recall.

Finally, example six from Table \ref{tableCosSimWordSubstitutionNouns} leads us to the paradigmatic relation between \textit{tabernacula} `tents' and \textit{nidos} `nests'. While a tent seems like a far-fetched analogy to a bird's nest and, consequently, has motivated various textual critics to assume an interpolation in the manuscript tradition, it is not our mission to solve such problems here. Instead, we rely on the text edition used in PROIEL and consider \textit{tabernacula} `tents' to be somewhat synonymous to \textit{nidos} `nests' in this very specific context. The underlying process is that of a metaphor with a source and target domain \parencite[275]{pantherMetaphorMetonymyLanguage2017}, probably much in the way as \textit{tabernaculum} `tent' is used to describe some people's commitment to philosophy\footnote{\url{https://cts.perseids.org/api/cts?request=GetPassage&urn=urn:cts:latinLit:phi0474.phi037.perseus-lat1:1.1-1.20}.}. If our intuition is correct and the analogy is restricted to a limited semantic use case, to be recognized by the limited range of words that can appear in the same sentence, then we should expect the two related target words to share barely any other contexts (or none at all). This hypothesis can be operationalized and tested using distributional measures like cosine distance in a \gls{CVS}.

Using the additional syntactic constraint of functional equivalence, we obtain, in addition to the results showed in Table \ref{tableCosSimWordSubstitutionNouns}, a list of further potential co-hyponyms (see Table \ref{tableAlmostIdenticalPROIELsentencesEquivalentSyntax}). From that new collection of results, we can quickly identify three cases (examples 1, 4 and 5) of near-synonymy. There is but a very narrow difference between the respective terms for bed, speech and inscription in these specific contexts. In example 4, we may note that, in the Latin WordNet, \textit{sermo} `talk' and \textit{verbum} `word' share a common hypernym synset\footnote{\url{https://latinwordnet.exeter.ac.uk/api/synsets/n/00019671/}.} at the third and fourth layer, respectively, of their hypernymy chain. Intuitively, we may think that their semantic distance (as indicated by their theoretical hypernymy chains) is rather moderate. However, as these text passages show us, the two can be almost indistinguishable in very specific contexts, such as being used in variations of a \textit{figura etymologica} where a major function of the accusative object is to intensify the verbal semantics \parencite[438]{mengeLehrbuchLateinischenSyntax2009}. Besides, we need to keep in mind that WordNet is a largely intuitive and theoretical resource, which may very well be contradictory to what we actually find in Latin language use. 

\begin{table}[ht]
	\begin{tabular}{p{\linewidth}}
    \ex[exno=PROIEL 62941]
        \begingl
            \gla toll-e \textbf{lect-um} tu-um et vad-e in dom-um tu-am //
            \glb lift-{\sc imp.sg} bed-{\sc n.sg.acc} your-{\sc n.sg.acc} and go-{\sc imp.sg} to house-{\sc f.sg.acc} your-{\sc f.sg.acc} //
            \glft take your bed and go to your house //
        \endgl
        \xe \\
	\ex[exno=PROIEL 47413]
        \begingl
            \gla toll-e \textbf{grabatt-um} tu-um et vad-e in dom-um tu-am //
            \glb lift-{\sc imp.sg} mattress-{\sc n.sg.acc} your-{\sc n.sg.acc} and go-{\sc imp.sg.} to house-{\sc f.sg.acc} your-{\sc f.sg.acc} //
            \glft take your mattress and go to your house //
        \endgl
    \xe \\ \hline
	\ex[exno=PROIEL 13284]
        \begingl
            \gla qui autem minor e-st in regn-o \textbf{cael-orum} maior e-st ill-o //
            \glb who but lesser be-{\sc 3sg} in kingdom-{\sc abl} heaven-{\sc gen.pl} greater be-{\sc 3sg} he-{\sc m.sg.abl} //
            \glft but who is lesser in the kingdom of heavens is greater than him //
        \endgl
    \xe \\
	\ex[exno=PROIEL 16820]
        \begingl
            \gla qui autem minor e-st in regn-o \textbf{De-i} maior e-st ill-o //
            \glb who but lesser be-{\sc 3sg} in kingdom-{\sc abl} god-{\sc gen} greater be-{\sc 3sg} he-{\sc m.sg.abl} //
            \glft but who is lesser in the kingdom of God is greater than him //
        \endgl
    \xe \\ \hline
    \ex[exno=PROIEL 13515]
        \begingl
            \gla et ai-t \textbf{puer-is} su-is //
            \glb and say-{\sc pf.3sg} servant-{\sc dat.pl} his-{\sc dat.pl} //
            \glft and he told his servants //
        \endgl
    \xe \\
	\ex[exno=PROIEL 11363]
        \begingl
            \gla et ai-t \textbf{discipul-is} su-is //
            \glb and say-{\sc pf.3sg} disciple-{\sc dat.pl} his-{\sc dat.pl} //
            \glft and he told his disciples //
        \endgl
    \xe \\ \hline
    \ex[exno=PROIEL 14018]
        \begingl
            \gla interroga-bo vos et ego un-um \textbf{sermo-nem} //
            \glb ask-{\sc fut.1sg} you also I one-{\sc m.sg.acc} talk-{\sc m.sg.acc} //
            \glft I also will ask you one thing //
        \endgl
    \xe \\
	\ex[exno=PROIEL 11133]
        \begingl
            \gla interroga-bo vos et ego un-um \textbf{verb-um} //
            \glb ask-{\sc fut.1sg} you also I one-{\sc n.sg.acc} phrase-{\sc n.sg.acc} //
            \glft I also will ask you one thing //
        \endgl
    \xe \\ \hline
    \ex[exno=PROIEL 14106]
        \begingl
            \gla cu-ius e-st imago h-aec et \textbf{suprascripti-o} //
            \glb who-{\sc m.sg.gen} be-{\sc 3sg} image this-{\sc f.sg.nom} and inscription-{\sc f.sg.nom} //
            \glft whose image and inscription is this? //
        \endgl
    \xe \\
	\ex[exno=PROIEL 11183]
    	\begingl
            \gla cu-ius e-st imago h-aec et \textbf{inscripti-o} //
            \glb who-{\sc m.sg.gen} be-{\sc 3sg} image this-{\sc f.sg.nom} and inscription-{\sc f.sg.nom} //
            \glft whose image and inscription is this? //
        \endgl
    \xe \\ \hline
    \ex[exno=PROIEL 14124]
    	\begingl
        	\gla responde-ns autem \textbf{Iesus} ai-t ill-is //
        	\glb answer-{\sc part.pres} but Jesus say-{\sc perf.3sg} he-{\sc m.dat.pl} //
          	\glft but Jesus told them in response //
        \endgl
    \xe \\
    \ex[exno=PROIEL 14548]
    	\begingl
        	\gla responde-ns autem \textbf{praeses} ai-t ill-is //
        	\glb answer-{\sc part.pres} but governor say-{\sc perf.3sg} he-{\sc m.dat.pl} //
          	\glft but the governer told them in response //
        \endgl
    \xe
	\end{tabular}
	\caption{Nouns in almost identical sentences from PROIEL, with equivalent syntactical structure. Target tokens are highlighted in bold.}
	\label{tableAlmostIdenticalPROIELsentencesEquivalentSyntax}
\end{table}

Generally, our choice of temporarily neglecting semantic distance for the consideration of co-hyponymy forces us to regard many (near-)synonyms, antonyms, etc., also as co-hyponyms. Consequently, we have to allow multiple relationships between two entities of our framework, namely multiple semantic relations between two entities (i.e., lexemes or, more empirically, tokens). Such a complexity is problematic in applied linguistics \parencite[2]{hendrickxSemeval2010Task82019} and machine learning \parencite[5836]{zellersNeuralMotifsScene2018} because it leads to a more ambitious classification target for computational models. One way to avoid and reduce that complexity is to consider each relation separately in a graph setting and combine all relation graphs into a large multigraph in the end \parencites{krauseSargraphsLanguageResource2016}{pinterPredictingSemanticRelations2018}, which is compatible with our abstract requirement of having multiple semantic relations between two entities.

Examples three and six from Table \ref{tableAlmostIdenticalPROIELsentencesEquivalentSyntax} are similar to the \textit{Iesus - scriba} `Jesus - scribe' case from Table \ref{tableCosSimWordSubstitutionNouns}: Both entities are animate and (at least somewhat) human, which is sufficient as a precondition to act as speakers in dialogical contexts. The second example, on the other hand, is a bit harder to tackle: Both \textit{caelorum} `of heavens' and \textit{Dei} `of god' are nominal modifiers to \textit{regno} `kingdom'. Since, in the New Testament, the Christian god is imagined as a ruler of heaven, we may conclude that the first instance points to the sphere or region of his rule, while the second one emphasizes the ruler himself. In contrast to most other examples shown here, they share a syntagmatic rather than a paradigmatic relation. We can observe this in the distribution of types on the sentence level:

\begin{itemize}
    \item \textit{caelorum}: ('\textbf{regnum}', 25), ('est', 21), ('in', 20), ('et', 19), ('qui', 14), ('simile', 9), ('\textbf{regno}', 7), ('autem', 7), ('homini', 6), ('enim', 5), ('vobis', 5), ('de', 4), ('quia', 4), ('non', 4), ('suo', 4), ('adpropinquavit', 3), ('dico', 3), ('nisi', 3), ('\textbf{intrabit}', 3), ('quod', 3), ('maior', 3), ('a', 3), ('agro', 3)
    \item \textit{Dei}: ('et', 537), ('in', 460), ('qui', 214), ('est', 208), ('non', 120), ('ut', 90), ('per', 86), ('fuit', 83), ('ad', 79), ('quae', 79), ('autem', 76), ('sunt', 70), ('quia', 67), ('cum', 65), ('\textbf{Dei}', 64), ('enim', 64), ('sed', 59), ('\textbf{Iesu}', 59), ('vobis', 56), ('\textbf{regnum}', 55), ('vos', 53), ('ne', 51), ('quod', 50), ('\textbf{Christi}', 50), ('de', 45), ('que', 42), ('quoniam', 39), ('si', 39), ('a', 38), ('ex', 33), ('ergo', 31), ('omnibus', 31), ('nos', 31), ('\textbf{Filius}', 31), ('\textbf{verbum}', 31), ('secundum', 31)
\end{itemize}

Both distributions are sorted by type frequency in descending order and cut off at a frequency of 2 and 30, respectively, for better legibility. In general, the frequencies are higher for \textit{Dei} `of god' because it occurs more often in the PROIEL corpus. This also affects the ranking of types (see \hyperref[functionWordsZipfDistribution]{above}), which is dominated in both cases by function words such as \textit{est} `he/she/it is', \textit{in} `in' or \textit{et} `and'. The case of \textit{caelorum} `of heavens' is a bit special, though, because the distribution ranking is lead by \textit{regnum} `kingdom', a noun that it typically modifies (`kingdom of heavens'). This prominent position in the ranking can be seen as an indication of idiomatization, i.e., the metaphor of the Christian god as a ruler and of heaven as the land which he rules has become the primary understanding and purpose of the type \textit{caelorum} `of heavens'. This is supported by the fact that morphological variants of the modified noun (such as \textit{regn-o} `kingdom-{\sc sg.abl}') follow not far behind in the ranking. While \textit{regnum} `kingdom' is also a frequent collocate for \textit{Dei} `of god', it does not lead the ranking and is accompanied by other equally frequent types like \textit{Iesu} `of Jesus', \textit{Christi} `of Christ' or \textit{Filius} `son'. Thus, the term's meaning is broader and more flexible, which correlates with its higher overall frequency in the corpus (see \hyperref[highFrequencyhighPolysemy]{below}).

Besides those general considerations, we may also note that both terms have strikingly distinctive collocates like \textit{intrabit} `he/she/it will enter' (sc. heaven) and \textit{verbum} `word' (sc. of god), which cannot be easily combined with the respective other term. Finally, \textit{Dei} `of god' collocates significantly with itself (burstiness), which can be a hint of ongoing definitional processes \parencite[361]{brookesCorrelationCollocationCohesion2020} and a theocentric perspective on religion \parencite[11]{callisterKeynessCollocationDistinctions2021}, both of which are plausible explanations in the case of the New Testament. Considering our focus on hyponymy, there is no intuitive way of discarding such examples from the search without sacrificing a lot of true positives. For the time being, it is worthwhile to accept a few hard negatives in exchange for increased precision and recall.

\section{Measuring Semantic Distance}
While we decided to consider, in principle, any semantic distance as eligible for hyponymy relations, we will nevertheless try to leverage measurements of such distance to improve our detection heuristics. One way of operationalizing semantic distance in a distributional framework lies in the comparison of lexical contexts in a syntactically confined space, such as segment-level token co-occurrences in a text corpus. This kind of approach will typically leave us with thousands of pairwise distributions to be assessed. To do so in a systematic manner, we need a suitable statistical methodology, in the form of \gls{DSM}s. One highly relevant \gls{NLP} task in that regard is language modeling. Traditional approaches have been around since at least 1980 \parencite[1270]{rosenfeldTwoDecadesStatistical2000} and, more recently, include combinations of pre-training and fine-tuning \parencite[1547]{hintonFastLearningAlgorithm2006}, which have been overwhelmingly popular since the emergence of ever larger language models \parencites{devlinBERTPretrainingDeep2019}{lewisBartDenoisingSequencetosequence2019}{songMassMaskedSequence2019}{ben-davidPERLPivotbasedDomain2020}{chungRethinkingEmbeddingCoupling2020}{joshiSpanBERTImprovingPretraining2020}{kaleNmT5IsParallelData2021}{luKELMKnowledgeEnhanced2021}{wangERNIE30Titan2021}{zouControllableGenerationPretrained2021}{cuiPERTPretrainingBERT2022}{zhangOPTOpenPretrained2022}. However, reusing pre-trained models is not always an option in corpus linguistics. Since we are interested in semantics as evident from the PROIEL corpus, we must not use models with external information because it would contaminate (and possibly override) the distributional information contained in our relatively small text collection. 
%TODO: replace OPT with an encoder decoder model like BERT, because its token embeddings are more suitable for hyponymy detection
Thus, we proceed by reusing the existing OPT model architecture \parencite{zhangOPTOpenPretrained2022}, but training from scratch instead of reusing random readily available parameters. 

Following recent standard practice, we provide a model card \parencite[222]{mitchellModelCardsModel2019} for our trained model, based on the official information provided by Meta for OPT and Llama\footnote{\url{https://huggingface.co/facebook/opt-350m}; \url{https://github.com/facebookresearch/llama/blob/main/MODEL_CARD.md}}:
\subsection{Model Card: Latin OPT}
\subsubsection{Model Details}
\begin{itemize}[font=\bfseries]
    \item[Model Developer] Konstantin Schulz
    \item[Variations] Latin OPT is a single model with 41 million parameters.
    \item[Input] Models input text only.
    \item[Output] Models generate text only. 
    \item[Model Architecture] Latin OPT is an auto-regressive language model that uses an optimized transformer architecture. \\
    \begin{table}[H]
        \centering
        \begin{tabular}{c|c|c|c|c}
        Training Data & Params & Content Length & Tokens & LR \\ \hline
        PROIEL treebank & 41M & 512 & 220k & 3.0 x $10^{-4}$
        \end{tabular}
        \caption{Model Architecture for Latin OPT, where `Params' is short for the overall number of model parameters and `LR' represents the initial learning rate.}
        \label{tab:latin_opt_architecture}
    \end{table}
    \item[Model Dates] Latin OPT was trained in August 2023.
    \item[Status] This is a static model trained on an offline dataset. No future versions are planned at the moment.
    \item[License] MIT license, available at: \url{https://github.com/facebookresearch/metaseq/blob/main/LICENSE}
    \item[Research Paper] More information can be found in the paper "OPT: Open Pre-trained Transformer Language Models", available at \url{https://arxiv.org/abs/2205.01068}.
    \item[Questions or comments] Instructions on how to provide feedback or comments on the model can be found in the \href{https://github.com/facebookresearch/metaseq#support}{support} section. 
\end{itemize}
\subsubsection{Intended Use}
Latin OPT can be used for prompting for evaluation of downstream tasks as well as text generation. In addition, the model can be fine-tuned on a downstream task.
\subsubsection{Hardware and Software}
\begin{itemize}[font=\bfseries]
    \item[Training Factors] We used Google Colab\footnote{\url{https://colab.research.google.com/}} for pretraining. Application and evaluation were performed on private hardware.
    \item[Carbon Footprint] Pretraining utilized a cumulative 4 GPU hours of computation on hardware of type T4-16GB (TDP of 70W). Estimated total emissions were 140 gCO2eq. 
\end{itemize}
\subsubsection{Training Data}
\begin{itemize}[font=\bfseries]
    \item[Overview] Latin OPT was trained on 220,000 tokens of data from the PROIEL treebank.
    \item[Data Freshness] The training data originates from the 2018 release of the PROIEL treebank. 
\end{itemize}
\subsubsection{Evaluation Results}
We did not perform any evaluation on public benchmarks because Latin OPT is too specialized for generic datasets.
\subsubsection{Ethical Considerations and Limitations}
Latin OPT is a new technology that carries risks with use. Testing conducted to date has been in Latin, and has not covered, nor could it cover all scenarios. For these reasons, as with all LMs, Latin OPT’s potential outputs cannot be predicted in advance, and the model may in some instances produce inaccurate, biased or other objectionable responses to user prompts. Therefore, before deploying any applications of Latin OPT, developers should perform safety testing and tuning tailored to their specific applications of the model.

\subsection{Cosine Similarity between Segment Pairs}
Using the newly trained model, we can compute the cosine distance between vector embeddings. Ideally, these embeddings should be calculated per token, not per segment. Unfortunately, our first approach at tokenization uses a vocabulary that is derived from Byte Pair Encoding \parencite[1716]{sennrichNeuralMachineTranslation2016}, i.e., an inventory of 10,000 byte sequences. This inventory starts out with a collection of all distinct bytes ($\sim$ characters) in the corpus and then continues to fill empty slots by merging the most frequent byte combinations into longer units. Since a vocabulary of 10,000 is rather large for a small text corpus like PROIEL (31,504 types), the tokenizer is able to represent most tokens using a single entry from the inventory. In some cases, frequent combinations of types (collocations etc.) get merged as well, leading to single representations for groups of tokens. As a consequence, the vectors that we receive as output by the model cannot be mapped consistently to single tokens anymore. Two possible countermeasures are 
\begin{enumerate}
    \item prohibiting merges of byte sequences, i.e., operating with a tiny vocabulary that only covers single characters, but not character sequences, or
    \item tokenizing only on the basis of entire types instead of bytes.
\end{enumerate}

Option 2 would work nicely in our confined setup, but fails to generalize to unseen data (out-of-vocabulary types; see \cite[57]{paradaSpokenTermDetection2010}). Besides, the vocabulary would have to grow infinitely corresponding to the size of training data, which is impractical as it would lead to much larger models, thereby increasing the time and cost of the calculations. Option 1, on the other hand, generalizes quite nicely at first glance: any text can be represented by the universal inventory of bytes. However, this also makes the model less expressive: Types in fusional languages such as Latin can become rather lengthy, like \textit{honorificentissimam} `most honorable-{\sc f.sg.acc}' (PROIEL 225208). Representing them byte by byte means that the final global representations for each byte will be informed by all occurrences of said byte, which probably amounts to many rather different types. Thus, the linguistic information contained in the byte representation will be far away from our traditional notions of characters, types or lemmata \parencite[1718]{sennrichNeuralMachineTranslation2016}. Ultimately, since we really want a stable deterministic detokenization, we need to retrain the (model and) tokenizer, forcing them to respect token boundaries (such as whitespace) in their merging rules.

Using the updated model and tokenizer, we obtain the cosine similarities for the text pairs provided in tables \ref{tableCosSimWordSubstitutionNouns} and \ref{tableAlmostIdenticalPROIELsentencesEquivalentSyntax} (see Table \ref{tab:cosSim1TokChangedProiel}). By sorting the similarity values in descending order, we get the direct visual impression of an inverted pyramid shape, which indicates a correlation of cosine similarity with text length, a well-known problem in research on semantic comparisons of longer texts \parencite[40]{singhalModernInformationRetrieval2001}. A quick statistical analysis reveals Pearson's correlation coefficient to be 0.57 for the cosine values and the pairwise mean character count, which can be interpreted as a solid positive correlation. Conceptually, this is not necessarily what we want: longer texts can be semantically more precise than shorter ones, thus not suggesting disregard for the purpose of hyponymy detection solely based on their length per se.

\begin{table}[ht]
    \centering
    % move entire table to the left, so all texts are displayed correctly
    \hspace*{-3cm}
    \begin{tabular}{c|c|c}
    Cosine Similarity & Sentence 1 & Sentence 2 \\ \hline
    0.99913 & vulpes foveas habent et volucres caeli \textbf{tabernacula} & vulpes foveas habent et volucres caeli \textbf{nidos} \\
    0.99909 & cuius est imago haec et \textbf{suprascriptio} & cuius est imago haec et \textbf{inscriptio} \\
    0.9967 & et ecce vox de \textbf{caelis} dicens & et ecce vox de \textbf{nube} dicens \\
    0.97526 & tolle \textbf{lectum} tuum et vade in domum tuam & tolle \textbf{grabattum} tuum et vade in domum tuam \\
    0.96394 & qui autem minor est in regno \textbf{caelorum} maior est illo & qui autem minor est in regno \textbf{Dei} maior est illo \\
    0.95523 & interrogabo vos et ego unum \textbf{sermonem} & interrogabo vos et ego unum \textbf{verbum} \\
    0.92421 & nonne et \textbf{publicani} hoc faciunt & nonne et \textbf{ethnici} hoc faciunt \\
    0.91422 & et ait \textbf{pueris} suis & et ait \textbf{discipulis} suis \\
    0.8273 & sexta \textbf{feria} in Syon & sexta \textbf{die} in Syon \\
    0.72061 & respondens autem \textbf{Iesus} ait illis & respondens autem \textbf{praeses} ait illis \\
    0.59476 & et ait illi \textbf{Iesus} & et ait illi \textbf{scriba} \\
    0.52413 & et respondens \textbf{centurio} ait & et respondens \textbf{eis} ait \\
    \end{tabular}
    \caption{Cosine similarities for selected sentence pairs from PROIEL, with target tokens highlighted in bold}
    \label{tab:cosSim1TokChangedProiel}
\end{table}

To verify this problematic correlation systematically at the corpus level, i.e., beyond the inspection of singular examples, we
\begin{enumerate}
    \item take the $\sim$ 30,000 types in the PROIEL corpus,
    \item discard all Hapax Legomena because we need at least two instances for comparison,
    \item for each type, extract all segments where this type occurs and order them by text length,
    \item calculate cosine similarity for neighboring segment pairs and specify their common text length to be the average of their individual text lengths,
    \item calculate the average cosine similarity for each text length (combined for all types; see Fig. \ref{figure:cosineSimilarityVsTextLength}),
\end{enumerate}

\begin{figure}[ht]
	\includegraphics[width=\linewidth]{cosine_similarity_vs_text_length.pdf}
	\caption{Correlation of cosine similarity and mean pairwise average text length. Each blue dot represents the average pairwise cosine similarity between segment pairs with a given average text length in characters. The red line indicates a linear regression of average cosine similarity and text length.}
	\label{figure:cosineSimilarityVsTextLength}
\end{figure}

Using this approach, we may observe at least two things:
\begin{enumerate}
    \item Cosine similarity tends to exhibit higher variance for longer texts. The dispersion of blue dots in the graph from left to right resembles a cone shape. This is possibly true for very short texts as well, but cannot be reasonably discerned from this mode of presentation, especially with regard to the scale of the x-axis.
    \item There is a positive correlation of cosine similarity and text length, even at the corpus level. $\sim$ 800 characters of additional text length seem to lead to an increase of 0.2 in cosine similarity. By definition, this correlation is limited to the point where all examples get a value of 1. An accumulation of this trend can be seen in the range of texts with length 400-800, where relatively many instances attain that highest value, making comparisons based on cosine similarity rather ineffective. Since the mean segment length (in characters) is 71 for the entire PROIEL corpus and we find relatively few segments of length >200 (see Fig. \ref{figure:SegmentLengthsProiel}), we may assume that this extreme manifestation is negligible for our purposes.
\end{enumerate}

\begin{figure}[ht]
	\includegraphics[width=\linewidth]{segment_lengths_proiel.pdf}
	\caption{Number of segments for various amount of text lengths. Text length is presented using a log scale to optimize the view of the distribution for the most frequent cases in PROIEL.}
	\label{figure:SegmentLengthsProiel}
\end{figure}

%TODO: distance between closest common hypernym for some of the examples, from Python "wordnet.py"
% TODO: continue here...

\section{High Variance in Low-Frequency Types: A Case Study on Dis Legomena}
In a first attempt to analyze possible confounding variables in our approach to Latin polysemy, we look at the occurrence frequencies for every lemma and make two important observations (see Fig. \ref{figure:cosineDistanceVsFrequency}):
\begin{itemize}
	\item Lemmata with a higher frequency of occurrence tend to exhibit a higher average cosine distance between the vectors of their usage contexts.
	\label{firstObservation}
	\item Lemmata with a lower frequency of occurrence have a higher variance, which leads to more outliers compared to the arithmetic mean. Conversely, highly frequent lemmata tend to be more predictable in their average cosine distance.
\end{itemize}
\begin{figure}[ht]
	\includegraphics[width=\linewidth]{cos_dist_vs_freq}
	\caption{Correlation of average cosine distance and frequency of occurrence. Each blue dot represents the average cosine distance between all vectors for one lemma. Each vector was computed from one usage context of the lemma in the corpus. The red line indicates a linear regression of average cosine distance and usage frequency.}
	\label{figure:cosineDistanceVsFrequency}
\end{figure}
The second observation may serve as a reminder that we have to be particularly cautious when looking at low-frequency lemmata: Their relatively high variance in average cosine distance indicates a widely varying degree of polysemy. We can validate our operationalization of polysemy through cosine distance by manually inspecting conspicuous cases: Lemmata that occur twice in the whole corpus and only within the same sentence are, by definition, the least polysemous in our setting. The reasons for this behavior are as follows:
\begin{itemize}
	\item Hapax Legomena have been discarded from the inspection because they only have one vector and cosine distance is not easily applicable to anything but pairs of vectors.
	\item The vectors are computed using distributional information (i.e. which other words appear in the sentence) and positional encoding. If a lemma appears twice in a sentence, the vectors for the two instances will differ mainly in their positional encoding, but not so much in the distributional information. This leads to very similar vectors with a low cosine distance.
\end{itemize}
Therefore, lemmata with 2 occurrences that both stem from the same sentence (e.g. \textit{Augustofratensis, Constantinopolis} or \textit{Sychem}, found in PROIEL 57662, 57780 and 23738, respectively) are prime examples for a low degree of polysemy. The phenomenon of low-frequency words tending to appear multiple times in quick succession (instead of a smooth distribution over the whole corpus) is referrd to as burstiness. Borrowing terminology from formal semantics, we may state that this burstiness is associated with a certain degree of logicality \parencite[178f.]{vonfintelFormalSemanticsGrammaticalization1995} or permutability \parencite[2]{altmannWordFrequencyBursts2009}: Words with a strong tendency towards burstiness are more logical and less permutable, i.e., more stable in their meaning even if the discourse context changes. \textit{Augustofratensis, Constantinopolis} and \textit{Sychem} are geographical references to ancient cities or provinces. Their meaning does not change, regardless of whether the respective entity is being founded, conquered, destroyed or acts as the birthplace of some famous person.

Besides, it is worth noting that burstiness often results from a topical focus, which in turn influences lexis \parencite[99]{pierrehumbertBurstinessVerbsDerived2012}. This topical focus also exists in PROIEL to some extent, as we can quickly determine from a simple calculation: 1,271 lemmata are Dis Legomena, i.e. they occur exactly twice in the corpus. 37 of them occur only within a single sentence. Thus, the observed probability of encountering the second instance of a Dis Legomenon in the same sentence as the first one is $37 / 1,271 = 0.029$. If, on the other hand, we follow the open-choice principle \parencite[109f.]{sinclairCorpusConcordanceCollocation1991} and naively assume an equal probability of occurrence for lemmata in each token slot, our expectation would look like this:
\begin{itemize}
	\item In the whole corpus, there are 19,335 sentences with 217,797 tokens, which in turn derive from 8,872 lemmata.
	\item Assuming open choice, we keep the length of each sentence and the token frequencies fixed, but shuffle the tokens randomly across all sentences. Now, the chance of a specific token occurring in a specific sentence is $1 / 217,797$ per slot, which gets multiplied by the sentence length (in terms of token count). For a sentence with 11 tokens (which is the average sentence length in PROIEL), the chance of a specific lemma occurring in that sentence is $11 / 217,797 = 5.1e-5$. The occurrence of a second token with the same lemma in that sentence is then $(11 - 1) / (217,797 - 1) = 4.6e-5$. 1 is subtracted from the sentence length because one other slot is already taken by the first occurrence of the lemma, and 1 is also subtracted from the token count because one of the tokens has already been placed. Using this pattern, we can run the following calculation for every sentence in the corpus:
	
	$\frac{sentence\_length * (sentence\_length - 1)}{token\_count * (token\_count - 1)}$
    \item Applying this formula to the observed average sentence length in PROIEL (11 tokens) and multiplying by sentence count (19,335), we get a theoretical probability of 4e-5 per Dis Legomenon: Thus, we would expect only 0.05 (i.e., basically none) of the 1,271 Dis Legomena to occur twice in the same sentence. The observed value is about 727 times higher.\footnote{Note that using the average sentence length is only a rough approximation. In fact, due to the high variance in sentence length (standard deviation of 9.4 tokens) and the associated inspection paradox \parencite[2]{palInspectionParadoxStochastic2022}, a random token tends to land in a longer-than-average sentence. To arrive at a more robust representation of that mechanism, we need to actually apply the given formula to every single sentence and consider its individual length, then sum the results for all sentences. This way, we get a probability of 0.10 in the end, shrinking the gap between observation and expectation to a factor of 364, which is about half as large as the rough approximation given above.} A permutation test corroborates this: when shuffling all tokens randomly across the existing sentences 1,000 times (thus preserving both lemma frequencies and sentence lengths), we never observe more than 3 such Dis Legomena in any given hypothetical corpus variant.
\end{itemize}

The immense deviation from our expectations leads to the conclusion that the open-choice principle is not a good model for lexis in sentences of this Latin corpus. We may even go further and argue that many statistical measures cannot be reasonably applied to language corpora because their mathematical preconditions are not met by the data \parencite[110]{shadrovaMeasuringCoselectionalConstraint2020}. In particular, the words in a corpus are not distributed evenly and their relative frequencies do not converge after a certain amount of text or time \parencite[5]{debowskiNaturalLanguagePerigraphic2018}.

%TODO: Alternativen zu TTR für lexikalische Dichte?
\label{highFrequencyhighPolysemy}
The \hyperref[firstObservation]{first observation}, on the other hand, suggests a positive correlation between polysemy and usage frequency: Words that we use more frequently tend to be more flexible in their meaning from a distributional point of view \parencite[576]{crossleyDevelopmentPolysemyFrequency2010}. We can verify that claim by subsampling a specific amount of usage contexts for every type in the corpus and calculating the \gls{TTR} \parencites[229]{kettunenCanTypetokenRatio2014}[1]{hamiltonDiachronicWordEmbeddings2016} for them:

\begin{enumerate}
	\item Select all types from PROIEL with an occurrence frequency $\geq$10.
	\item For every selected type, collect all tokens that occurred in the same sentences as that type.
	\item Omit all types did not occur together with at least 100 tokens.
	\item For the remaining types, randomly sample 100 tokens from their usage contexts.
	\item From the selected tokens, calculate the \gls{TTR} for each type.
\end{enumerate}

The thresholds of 10 and 100 were arbitrarily chosen. Their main purpose is to ensure that we have enough (i.e., 10) instances of a single type to avoid the large variance that became apparent in Fig. \ref{figure:cosineDistanceVsFrequency} and to mitigate the influence of text length effects (max. 100 tokens per type), which is a known problem for \gls{TTR} and other measures of lexical diversity \parencite[66]{koizumiRelationshipsTextLength2012}. Nevertheless, in the lower left corner of Fig. \ref{figure:ttr_vs_frequency} (\gls{TTR} <0.6, frequency <4), we can still observe some traces of higher lexical variance in low-frequency types. Generally, as indicated by the linear regression, high-frequency types tend to exhibit a slightly larger lexical diversity in their usage contexts. This underlines the importance of explicitly dealing with issues of polysemy in the extraction of hyponyms for language learning, as was \hyperref[high_frequency_hyponyms_for_language_learning]{outlined above}.

\begin{figure}[ht]
	\includegraphics[width=\linewidth]{ttr_vs_freq.pdf}
	\caption{Correlation of \gls{TTR} and frequency of occurrence. Each blue dot represents a specific type. \gls{TTR}s were calculated only for types that occur at least 10 times in the corpus and have at least 100 tokens in their aggregated usage contexts, i.e., sentences. For the chosen types, a random sample of 100 tokens was selected for the calculation of \gls{TTR} to avoid text length effects. The red line represents a linear regression.}
	\label{figure:ttr_vs_frequency}
\end{figure}

\section{Clustering of Contextual Embeddings}
As we outlined before, the actual clustering of usage contexts for every token in the PROIEL corpus will be performed on the lexeme level. The procedure is as follows:
\begin{enumerate}
	\item For every token in the corpus, compute its embedding vector in the sentence context.
	\item Initially, treat every embedding vector as its own cluster.
	\item For every pair of clusters, compute the similarity between their vectors. If a cluster contains more than one vector, compute the average for all of its vectors and use that for comparison to the other cluster.
	\item Merge those clusters that are very close to each other.
	\item Repeat steps 3 and 4 until a specific threshold is reached.
\end{enumerate}

Ideally, the recursive iterations should take place after every merging of two clusters. This way, the newly formed cluster has a chance to be merged again immediately (if its similarity to another cluster is still above the threshold). However, for computational reasons, we instead use small batches of 15 clusters and compute the similarity values between all distinct pairs for them. This leads to about 15,000 batches being simultaneously processed in every iteration. Thus, the clustering method loses some of its accuracy while being substantially faster to compute.

The initial number of clusters, where each usage context represents its own cluster, is 217,923 , which corresponds closely to the number of tokens in the PROIEL corpus (217,797). 
%TODO: explain the difference (126 or 0.06%) between usage contexts / clusters and PROIEL tokens
By using a threshold of 0.9 for the cosine similarity that is required to merge two clusters, exactly 1 new cluster emerges. Correspondingly, the overwhelming majority of clusters was not merged at all. The merged cluster contains two usage contexts of the word \textit{et} `and':

\vspace{10pt}
\ex[exno=PROIEL 11203]
	\begingl
		\gla et responde-ns Iesus ai-t ill-is //
		\glb and reply-{\sc part.pres} Jesus say.{\sc impf-3sg} that-{\sc pl.dat/abl} //
		\glft and Jesus replied to them by saying //
	\endgl
\xe 
\vspace{10pt}
\ex[exno=PROIEL 16513]
	\begingl
		\gla et responde-ns Iesus ai-t ill-i //
		\glb and reply-{\sc part.pres} Jesus say.{\sc impf-3sg} that-{\sc dat} //
		\glft and Jesus replied to him by saying //
	\endgl
\xe
\vspace{10pt}

The two sentences differ only in their lexis, and only in the last token. Even for those last tokens, the relevant lemma (\textit{ille} `he, that') is the same. Thus, the representations are so close to each other (cf. Table \ref{tableCosSimWordSubstitution}) that they get merged into a larger cluster quite easily.

This changes when we set the similarity threshold to 0.8 . To achieve a higher performance of the clustering algorithm, we do not perform this additional processing step on the original clusters that were all unmerged. Instead, we use the results of the first clustering phase, which is not important in this early phase (because it barely differs from the original clusters), but will be crucial later on. This procedure yields 41 merged clusters, some of which contain up to 4 usage contexts. All in all, the merged vectors now include 88 (0.04\%) usage contexts. Interestingly, the merges occurred not just for further instances of \textit{et} `and' (e.g., PROIEL 10565 and 48782) or other function words like \textit{ibi} `there' (PROIEL 57777 and 78440), but also for seemingly unrelated pairs like \textit{videam -- diligentiam} `I may see -- carefulness': 

\vspace{10pt}
\ex[exno=PROIEL 77240]
	\par
	\begingl
		\gla itaque exspect-o Thessalonic-ae act-a Kal Sext ex qu-ibus statu-a-m in tu-os ne agr-os confugi-a-m ut ne que vide-a-m homin-es qu-os nol-i-m et te ut scrib-is vide-a-m et prop-ius si-m si quid ag-a-tur id que intelle-x-i cum t-ibi tum Quint-o fratr-i placere an abe-a-m Cyzic-um //
		\glb thus wait-{\sc 1sg} Thessaloniki-{\sc loc} transaction-{\sc pl} first.day august from which-{\sc pl.abl} decide-{\sc sbjv-1sg} into your-{\sc pl.acc.m} whether field-{\sc acc.pl} flee-{\sc sbjv-1sg} so.that not =and see-{\sc sbjv-1sg} man-{\sc pl} who-{\sc pl.acc.m} be.unwilling-{\sc sbjv-1sg} and you.{\sc acc} as write-{\sc 2sg} see-{\sc sbjv-1sg} and near-{\sc comp} be.{\sc sbjv-1sg} if something conduct-{\sc sbjv-3sg.pass} this.{\sc n} =and  perceive-{\sc perf-1sg} when you-{\sc dat} then Quintus-{\sc dat/abl} brother-{\sc dat} please or go.away-{\sc sbjv-1sg} Cyzicus-{\sc acc} //
		\glft Thus, I wait at Thessaloniki for the transactions of August 1. On the basis of them, I will decide whether to take refuge at your estate, so that I may not have to see those that I do not want to, and instead see you (as you wrote to me), and be closer if anything happens (and this, I have learned, is pleasant for you and especially for Quintus), or to go away to Cyzicus //
	\endgl
\xe 
\vspace{10pt}

\ex[exno=PROIEL 77314]
	\par
	\begingl
		\gla in qu-o si iam nostr-a salus cum hac leg-e despera-t-a eri-t vel-i-m pro tu-o in me amor-e hanc inan-em me-am diligentia-m miserabil-em potius quam inept-am put-e-s sin es-t aliquid spe-i d-e-s opera-m ut ma-ior-e diligenti-a posthac a nostr-is magistrat-ibus defend-a-mur //
		\glb in which-{\sc abl.n/m} if already our-{\sc f} salvation with this.{\sc f.abl} law-{\sc abl} give.up-{\sc part.perf.pass-f.sg/n.pl} be.{\sc fut-3sg} want-{\sc sbjv-1sg} for.the.sake.of your-{\sc dat/abl.m/n} towards I.{\sc acc} love-{\sc abl} this.{\sc f.acc} useless-{\sc acc.m/f} my-{\sc f.acc} diligence-{\sc acc} pitiful-{\sc acc.m/f} rather than silly-{\sc f.acc} think-{\sc sbjv-2sg} if.however be-{\sc 3sg} anything hope-{\sc gen/dat} give-{\sc sbjv-2sg} effort-{\sc acc} so.that large-{\sc comp-abl} diligence henceforth by our-{\sc pl.dat/abl} official-{\sc pl.dat/abl} defend-{\sc sbjv-1pl.pass} //
		\glft Regarding this matter, if my salvation will already be hopeless together with this law, I would like you to think, because of your love towards me, that my useless diligence is pitiful rather than silly; if, however, there is any kind of hope, please try to make sure that, in the future, we will be defended by our officials with greater diligence //
	\endgl
\xe
\vspace{10pt}

The passages belong to Cicero's Letter to Atticus, book 3, letters 15 and 23. They share almost no vocabulary except for function words like \textit{tuus} `your', \textit{ut} `that, so that', or \textit{si} `if'. Their semantic similarity is possibly related to the overlap of involved persons (Cicero and Atticus), but Cicero's brother Quintus is mentioned in the first passage. Another factor is the narrative structure: both passages outline two potential future scenarios. In the first letter, Cicero ponders on his next action:
\begin{enumerate}[label=(\alph*)]
	\item go to Atticus' estate
	\item go to Cyzicus
\end{enumerate}
Correspondingly, the second letter shows two potential outcomes that Cicero wishes for:
\begin{enumerate}[label=(\alph*)]
	\item Atticus should feel pity for him (if all hope is relinquished)
	\item Atticus should get the officials on Cicero's side (if he gets the chance)
\end{enumerate}

The distribution of cluster size seems to follow a Zipfian distribution (see Fig. \ref{figure:cluster_size_distribution}). 

\begin{figure}[ht]
	\includegraphics[width=\linewidth]{cluster_size_distribution.pdf}
	\caption{Distribution of cluster size for merged clusters at a similarity threshold of 0.9 . The largest cluster (128 usage contexts) was excluded for better legibility.}
	\label{figure:cluster_size_distribution}
\end{figure}

All merged clusters consist of usage contexts for the same lemma, i.e., there are no cross-lexeme clusters. This does not change if we lower the similarity threshold to 0.8, 

%TODO: finish this sentence/paragraph
%TODO: mbert demo, document-level domain embeddings, glossing; explain custom glosses?

%As for the gold standard on hyponymy, the Latin part of the WordNet was used. In a very basic approach, we will extract from it the explicit hyponym and hypernym pairs only for nouns, converting both to the notation of "hyponym,hypernym". This ignores additional properties like transitivity in order to create an explicit empirical baseline for later comparison. Besides, since there is no previous empirical knowledge from applying machine learning to hyponymy detection in Latin, we shall transform the complex problem to a much simpler one and then, gradually, reintroduce the complexity: At the very beginning, the machine's task will be to detect whether there are instances of hyponymy in a given sentence or not, which is called binary classification because there are only two options (either there are some instances or there are none). Eventually, the most complex question may look something like this: Given a Latin text passage, which words in it are hyponyms of which other ones and which word sense do they express in that specific context? Notably, the main dimensions of complexity are the following:
%\begin{itemize}
   %      \item The input can be a sentence (easy) or some longer kind of segment, e.g. a paragraph (difficult).
      %   \item The detection can be binary (easy) or token-based (difficult).
         %\item Polysemy can be leveled (easy) or explicitly disambiguated (difficult).
%\end{itemize}
%Furthermore, there are various ways of balancing to consider. By comparing plain text to the WordNet lemmata in a superficial manner, we can only detect hyponymy for those sentences in which the hyponyms appear in their base form. For languages with little inflectional variety (like English), this may be an acceptable choice. For Latin, however, with its rich morphology, such a simplification has far-reaching consequences: About 3\% of the sentences were associated with hyponymy (see Table ~\ref{table:hyponymyPercentageShallow}).

%\begin{table}[ht]
%	\begin{tabular}{ c | c | c }
  %		 & All Sentences & With Hyponyms \\ \hline
  %		Absolute & 9 992 000 & 304 677 \\
  %		Percentage & 100\% & 3.05\%  \\
	%\end{tabular}
	%\caption{Percentage of sentences with hyponymy, extracted by shallow comparison to WordNet}
	%\label{table:hyponymyPercentageShallow}
%\end{table}

%Therefore, machines will achieve high accuracy rates (>95\%) using rather simple rules, e.g. by never predicting a hyponym. Alternatively, they may produce similar results by memorizing all of the (most frequent) hyponyms: There are 6781 noun lemmata in WordNet that also appear as surface forms in the Corpus Corporum. As with many linguistic frequency distributions, their occurrence counts follow Zipf's law \parencite[VI]{zipfPsychobiologyLanguageIntroduction1936}. In other words, a few "super-hypernyms" like \textit{causa} (cause, reason), \textit{genus} (race, kind) or \textit{homo} (human being, man) dominate the overall picture, leaving rarely occurring nouns such as  \textit{concussio} (shaking, concussion),  \textit{audientia} (hearing, listening) or  \textit{congelatio} (freezing, congealing) far behind. This information will be helpful for language learning purposes because, in some cases, rather superordinate terms may be acquired later than their prototypical hyponyms \parencite[27]{ellisUsagebasedLanguageInvestigating2013}. This is true for the acquisition of both L1 and L2. In Latin, however, we are not concerned with acquisition, but with learning: Students usually do not speak Latin in their everyday life or during  classes, except when occasionally reading aloud a text. Reciting Latin texts does not (yet) seem to be a well-defined concept in terms of language learning \parencite[243, fn. 89]{munznerNeueWegeIm2013}.

%(see Table ~\ref{table:modelAccuracyShallow})
%\begin{table}[ht]
%	\begin{tabular}{ c | c | c | c }
  %		 & Precision & Recall & F1 \\ \hline
  %		Score & .73 & .41 & .53 \\
	%\end{tabular}
	%\caption{Evaluation of the simplest model, trained with data from a shallow comparison to WordNet}
	%\label{table:modelAccuracyShallow}
%\end{table}
%\begin{figure}[ht]
%\includegraphics[width=\linewidth]{freq_dist_hyponyms_cc_plain}
%\caption{Frequency distribution of hyponym lemmata in the Corpus Corporum (plain text)}
%\label{figure:frequencyDistributionHyponymsCCplain}
%\end{figure}
